% Discussion subsection: Visual remapping and the emergence of local retinotopy
% For LearningFlexibleActionControl paper
% ~800-1200 words

\subsection{Local retinotopy in non-retinotopic maps: a caveat for visual remapping}

A key implication of our framework concerns the interpretation of retinotopic
signals in cortical areas that are not primarily organized around spatial position.
Visual remapping---the predictive updating of receptive fields in anticipation of
saccades---was first described in LIP neurons \citep{duhamel1992} and has since been
observed across a broad network of areas including the frontal eye field (FEF)
\citep{umeno1997}, superior colliculus \citep{walker1995}, area V4
\citep{neupane2016}, and even early visual cortex \citep{nakamura2002}. The standard
interpretation of remapping relies on an assumption that is rarely made explicit:
that the areas exhibiting remapping maintain a well-defined retinotopic map whose
coordinates can be predictively updated by a corollary discharge signal
\citep{sommer2006, marino2021}. However, recent evidence of pervasive retinotopic
organization in areas traditionally considered non-retinotopic---including
inferotemporal cortex \citep{arcaro2017}, scene-selective regions, and face patches
\citep{arcaro2017}---raises the question of what \emph{kind} of spatial organization
is needed to support remapping-like phenomena.

Our topological alignment framework, together with a complementary computational
analysis, suggests that this assumption deserves careful scrutiny. In a separate
analysis, we trained a standard self-organizing map (SOM) on a dataset of colored
rectangles varying along multiple dimensions---shape, orientation, size, aspect
ratio, position, and color. The map was thus driven to organize itself around the
manifold of shape features, not around spatial position per se. Yet, when probed
with localized stimuli (Gaussian peaks at known spatial positions), the trained
network revealed statistically overwhelming spatial smoothness: nearby units in the
grid responded to stimuli at similar retinal positions (Spearman $\rho \approx 0.5$,
$Z > 17$), despite the fact that the absolute range of receptive field centers was
tiny---a fraction of the image extent. In other words, the map exhibited
\emph{local} retinotopic structure without being \emph{globally} retinotopic.

This finding has a straightforward geometric explanation. A self-organizing map
preserves the topology of its training manifold: nearby units encode nearby regions
of the input space. When the input data has spatial structure---as any visual input
inevitably does, since objects occupy positions on the retina---spatial position
becomes one embedded dimension among many in the learned manifold. The topographic
constraint of the SOM ensures that this spatial dimension is represented smoothly
across the grid, even when other dimensions (shape, color) dominate the
organization. Local retinotopy is thus not a dedicated map of space; it is a
\emph{geometric byproduct} of topographic learning on spatially structured input.

This observation has direct implications for how we interpret retinotopic signals
in the parietal and frontal cortex. As one ascends the visual hierarchy into
posterior parietal cortex (IPS0--IPS4) and beyond, receptive fields become
progressively larger and more feature-selective
\citep{swisher2007, arcaro2020}. Retinotopic organization
in these regions is coarser than in V1, with maps that are increasingly dominated by
task-relevant features such as attention priority, motor intention, and object
identity. This is precisely the regime our SOM analysis models: a map that has
organized itself around non-spatial features but retains a residual spatial signal
that is weak in absolute terms yet robust when assessed at the population level.

The caveat for visual remapping is this: the finding of spatially organized
responses in an area does not, by itself, demonstrate that the area maintains a
strict retinotopic coordinate system that is updated by a saccadic corollary
discharge. An alternative possibility is that the spatial signal reflects residual
local retinotopy within a map organized around other features---and that the
``remapping'' observed in such areas is better understood as a consequence of
topographic continuity rather than as evidence for a dedicated spatial coordinate
transformation. In our SOM analysis, the spatial signal was undetectable with
hard winner-take-all readout (activating only a handful of units) but became
strongly significant with soft population readout. This methodological point
parallels a well-known challenge in neurophysiology: single-neuron recordings in
parietal cortex may fail to reveal spatial selectivity because individual neurons
have large receptive fields dominated by non-spatial features, whereas
population-level analyses can recover the spatial structure hidden in the
distributed response \citep{sereno2001}.

This perspective connects directly to the topological alignment mechanism presented
in this paper. Our model's visual map---a supervised topological map (STM)
organizing retinal inputs within a shared internal space---is not built to be
retinotopic. It is built to align with somatosensory, proprioceptive, and policy
maps so that a common point in the internal space encodes a coherent sensorimotor
contingency. Yet, precisely because the visual input is spatially structured,
the visual STM will inevitably retain local spatial organization. Far from being
a nuisance or a vestige, this residual retinotopy is what makes cross-modal
alignment possible: if a visual map organized around object features completely
lost spatial information, it could not be aligned with the inherently spatial
maps of proprioception and touch. The residual spatial scaffold is a
\emph{necessary condition} for the alignment to succeed.

In summary, we suggest that the presence of retinotopic-like signals in areas
beyond early visual cortex should not be taken as evidence that these areas
function as spatial maps in the same sense as V1. Rather, local retinotopy may
be an intrinsic geometric property of any topographic representation trained
on spatially structured visual input. This does not diminish the functional
significance of these signals---on the contrary, it suggests that residual
spatial organization serves as the geometric scaffold enabling the cross-modal
alignment that our framework identifies as the basis for goal-directed action
control. Parietal areas are not retinotopic \emph{despite} processing shapes
and actions; they retain spatial structure \emph{because} the topology of visual
input, from which shapes and actions are learned, is itself spatial.

% References to add:
% \bibitem{duhamel1992} Duhamel, Colby, Goldberg (1992). Science.
% \bibitem{umeno1997} Umeno, Goldberg (1997). J Neurophysiol.
% \bibitem{walker1995} Walker, Fitzgibbon, Goldberg (1995). J Neurophysiol.
% \bibitem{neupane2016} Neupane, Guitton, Pack (2016). PNAS.
% \bibitem{nakamura2002} Nakamura, Colby (2002). PNAS.
% \bibitem{sommer2006} Sommer, Wurtz (2006). Brain Res Rev.
% \bibitem{marino2021} Marino, Mazer (2021). Annu Rev Vis Sci.
% \bibitem{arcaro2017} Arcaro, Livingstone (2017). J Neurosci.
% \bibitem{arcaro2020} Arcaro, Honey, Mruczek, Kastner, Hasson (2015). Cereb Cortex.
% \bibitem{swisher2007} Swisher, Halko, Merabet, McMains, Somers (2007). J Neurosci.
% \bibitem{sereno2001} Sereno, Pitzalis, Martinez (2001). Science.
